\documentclass[paper=a4, parskip=half-, ngerman, fontsize=11pt]{scrreprt}
\usepackage[T1]{fontenc}

\usepackage[ngerman]{babel}
\babelprovide[hyphenrules=ngerman-x-latest]{ngerman}

\usepackage{csquotes}
\usepackage[backend=biber]{biblatex}
\addbibresource{Proseminar.bib}

\usepackage{graphicx}
\usepackage{url}
\usepackage{amsmath}
\usepackage{booktabs}
\usepackage{xcolor}


\begin{document}

% Titelseite
\begin{titlepage}
    \begin{center}
        % Unilogo oben
        \includegraphics[width=0.5\textwidth]{logo_fernuni_hagen.png}\\[2cm]

        {\LARGE \textbf{Übertragung von Signalen auf elektrischen Leitungen}}\\[2cm]

        \textbf{Proseminar Mathematik in der Technik}\\
        \textbf{Modulnummer:} 61711\\
        Fakultät für Mathematik + Informatik\\[0.5cm]

        \begin{tabbing}
            \hspace{6cm} \= \kill
            \textbf{Name:} \> Sven Schmidt \\
            \textbf{Matrikelnummer:} \> 4125169 \\
            \textbf{Abgabedatum:} \> \today \\
            \textbf{Prüfer:} \> PD Dr.-Ing. Stefan Helfert \\
        \end{tabbing}

        \vfill

        {\large FernUniversität in Hagen}\\
        {\large Sommersemester 2025}
    \end{center}
\end{titlepage}

\chapter{Einführung}

\chapter{Modell des Zweidrahtleiters}

In dieser Ausarbeitung beschränken wir uns auf die Untersuchung der Ausbreitung elektromagnetischer Wellen in
Zweidrahtleitern (Paralleldraht), die schematisch in der folgenden Abbildung gezeigt.
\begin{figure}[!h]
    \begin{center}
        \includegraphics[width=0.5\textwidth]{images/Leiter.png}
        \caption{Elektrische Leiter mit ihren Anwendungen. Entnommen aus~\cite{FernuniSkript}.}
        \label{Leiter}
    \end{center}
\end{figure}

Aus den Gesetzen der Elektro- und Magnetostatik ist bekannt, dass sich beim Durchfließen eines elektrischen Stroms um
den Leiter ein magnetisches Feld ausbildet. Analog bildet die Potenzialdifferenz beider Leiter ein elektrisches Feld
aus. Die folgende Abbildung~\ref{Felder} zeigt den Querschnitt eines Leiters mit eingezeichneten Feldlinien.

\begin{figure}[!h]
    \begin{center}
        \includegraphics[width=0.5\textwidth]{images/Felder.png}
        \caption{Elektrische Leiter mit ihren Anwendungen. Entnommen aus~\cite{LeitungenUndFilter}.}
        \label{Felder}
    \end{center}
\end{figure}

Im Folgenden nehmen wir an, dass der Abstand $d$ zwischen den Leitern viel kleiner ist, als die Wellenlänge $\lambda$
der auftretenden Signale,
\[ d \ll \lambda = \frac{c_{0}}{f \sqrt{\epsilon_{r}}} \]
wobei $c_{0}$ die Lichtgeschwindigkeit im Vakuum, $f$ die Frequenz des Signals und $\epsilon_{r}$ die relative
Dielektrizitätskonstante bezeichnet. Wir machen hier die vereinfachende Annahme, dass $\epsilon_{r}$ eine (konstante)
Eigenschaft des Mediums ist, in dem sich das Signal ausbreitet. Im Vakuum gilt $\epsilon_{r} = 1$.

Wir betrachten nun ein kleines, infinitesimales Leitungsstück $\mathrm{d}s$,
\begin{figure}[!h]
    \begin{center}
        \includegraphics[width=0.5\textwidth]{images/Leitung1.png}
        \caption{Elektrische Leiter mit ihren Anwendungen. Entnommen aus~\cite{LeitungenUndFilter}.}
        \label{Leitung1}
    \end{center}
\end{figure}

Der Fluss $\Phi$ des magnetischen Feldes $\textbf{H}$ bewirkt eine Induktivität $L_{s}$
\[ \Phi_{s} = i \cdot L_{s} \], wobei $i$ der Strom ist.
Das elektrische Feld $\textbf{E}$ induziert Ladungen auf der Oberfläche des Leiters und damit eine Kapazität $C_{s}$,
die über die Spannung $u$ wie folgt zusammenhängt,
\[ Q_{s} = C_{s} \cdot u . \]
Die gesamte Leitung stellen wir uns aus diesen Leitungsstücken zusammengesetzt vor wie auf Abbildung~\ref{Leitung2}
dargestellt.
\begin{figure}[!h]
    \begin{center}
        \includegraphics[width=0.5\textwidth]{images/Leitung2.png}
        \caption{Elektrische Leiter mit ihren Anwendungen. Entnommen aus~\cite{LeitungenUndFilter}.}
        \label{Leitung2}
    \end{center}
\end{figure}
Die führt auf die folgenden qualitativen Eigenschaften: Wir eine Spannung am Leitungsanfang angelegt, muss zunächst die
erste Kapazität $Q_{s_{1}}$ aufgeladen werden. Erst dann bildet sich eine Spannung an der Längsinduktivität $L_{s_{1}}$
an. Der aufgebaute Strom wird nun $Q{s_{2}}$ aufladen etc., bis zum Leiterende. Die Ausbreitung des Einschaltvorgangs
findet mit endlicher Geschwindigkeit statt.

Ein elektrischer Strom, der durch einen Leiter fließt, erfährt einen Widerstand $R_{s}$ auf einer Strecke
$\mathrm{d}s$. Wenn wir annehmen, dass das den Leiter umgebende Dielektrikum leitend ist, dann bezeichnen wir dies mit
dem Querleitwert
$G_{s}$ auf der Strecke $\mathrm{d}s$. Die Abbildung~\ref{Leitung3}
\begin{figure}[!h]
    \begin{center}
        \includegraphics[width=0.5\textwidth]{images/Leitungsstueck.png}
        \caption{Elektrische Leiter mit ihren Anwendungen. Entnommen aus~\cite{LeitungenUndFilter}.}
        \label{Leitung3}
    \end{center}
\end{figure}



\chapter{Die Telegraphenleitungen}
Abbildung~\ref{Leitung4} zeigt nochmals ein infinitesimales Leitungsstück $\mathrm{d}s$ eines verlustbehafteten
Leiters. Am Leitungseingang liegen Strom $i$ und Spannung $u$ an, am Leitungsende $i + \mathrm{d}i$ und $u +
\mathrm{d}u$.
\begin{figure}[!h]
    \begin{center}
        \includegraphics[width=0.5\textwidth]{images/Leitungsstueck2.png}
        \caption{Einschaltbild eines infinitesimalen Leitungsstücks einer verlustbehafteten Leitung. Entnommen
        aus~\cite{LeitungenUndFilter}.}
        \label{Leitung4}
    \end{center}
\end{figure}

Wir definieren auf einem Stück $\mathrm{d}s$

\begin{tabular}{l l}
    den Widerstandsbelag & $R^{\prime} = \frac{R_{s}}{s}$ \\
\addlinespace
    den Induktivitätsbelag & $L^{\prime} = \frac{L_{s}}{s}$ \\
\addlinespace
    den Kapazitätsbelag & $C^{\prime} = \frac{C_{s}}{s}$ \\
\addlinespace
    und den Leitungsbelag & $G^{\prime} = \frac{G_{s}}{s}$. \\
\end{tabular}

Der Widerstand $R^{\prime}$ und die Induktivität $L^{\prime}$ verursachen einen Spannungsabfall von
\[
i \, R^{\prime} \, \mathrm{d}z + \frac{\partial i}{\partial t} \, L^{\prime} \, \mathrm{d} z
\]
entlang von $\mathrm{d}s$. Mit der Kirchhoffschen Maschengleichung {\color{red}REF} folgt für die Spannung
\begin{equation}
u(z, t) = R^{\prime} \, i(z, t) + L^{\prime} \, \frac{\partial i(z, t)}{\partial t} + u(z + \Delta z, t).
\end{equation}
Durch Querleitwert $G_{s}$ und Querkapazität $C_{s}$ fließt der Strom
\[
u \, G^{\prime} \, \mathrm{d}z + \frac{\partial u}{\partial t} \, C^{\prime} \, \mathrm{d} z.
\]
Mit der Kirchhoffschen Knotenregel {\color{red}REF} folgt
\begin{equation}
    i(z, t) = G^{\prime} \, u(z, t) + C^{\prime} \, \frac{\partial u(z + \Delta z, t)}{\partial t} + i(z + \Delta z, t).
\end{equation}
Wir formen nun beide Gleichungen um,
\begin{align}
\frac{u(z + \Delta z, t) - u(z, t)}{\Delta z} &= \frac{R^{\prime}}{\Delta z} \, i(z, t) + \frac{L^{\prime}}{\Delta z}
\,
\frac{\partial i(z, t)}{\partial t} \\
\frac{i(z + \Delta z, t) - i(z, t)}{\Delta z} &= \frac{G^{\prime}}{\Delta z} \, u(z, t) + \frac{C^{\prime}}{\Delta z}
\, \frac{\partial u(z + \Delta z, t)}{\partial t}.
\end{align}
Mit
\[
\lim_{\Delta x \to 0} \frac{u(x+\Delta x, t)-u(x, t)}{\Delta x} = \frac{\partial u(x, t)}{\partial x} \quad \text{und}
\quad \lim_{\Delta x \to 0} \frac{\Delta R}{\Delta x} = R' \quad \text{folgt}
\]
\begin{align}
    \frac{\partial u}{\partial z} &= -\left(R^{\prime} + L^{\prime}\frac{\partial}{\partial t}\right)i \label{eq:Dgl1}
    \\
    \frac{\partial i}{\partial z} &= -\left(G^{\prime} + C^{\prime}\frac{\partial}{\partial t}\right)u \label{eq:Dgl2}
\end{align}
Dies ist ein System partieller linearer Differentialgleichungen 1.~Ordnung mit konstanten Koeffizienten. Sie werden
auch Differentialgleichungen der elektrischen Leitung genannt.
Gleichung~\eqref{eq:Dgl1} beschreibt den Spannungsabfall verursacht durch den Leitungswiderstand $R^{\prime}$ und
Induktivität $L^{\prime}$, siehe Abbildung~\ref{Leitung4}.
Die zweite Gleichung~\eqref{eq:Dgl2} beschreibt den Stromfluss zwischen den Leitern, verursacht durch Kapazität
$C^{\prime}$ und Querleitwert $G^{\prime}$.

Oftmals werden diese Gleichungen in die sogenannten Telegraphenleitungen weiter umgeformt. Dazu leiten wir beide
Gleichungen einmal nach $x$ und einmal nach $t$ ab,
\begin{align}
    \frac{\partial^{2} u}{\partial x^2} &= -R^{\prime} \, \frac{\partial i}{\partial x} - L^{\prime} \,
    \frac{\partial^2 i}{\partial t \,
    \partial x} \label{eq:Dgl3} \\[1.5ex]
    \frac{\partial^{2} i}{\partial x^2} &= -G^{\prime} \, \frac{\partial u}{\partial x} - C^{\prime} \,
    \frac{\partial^2 u}{\partial t \,
    \partial x} \label{eq:Dgl4} \\[1.5ex]
    \frac{\partial^{2} u}{\partial x \, \partial t} &= -R^{\prime} \, \frac{\partial i}{\partial t} - L^{\prime} \,
    \frac{\partial^2
    i}{\partial t^{2}} \label{eq:Dgl5} \\[1.5ex]
    \frac{\partial^{2} i}{\partial x \, \partial t} &= -G^{\prime} \, \frac{\partial u}{\partial t} - C^{\prime} \,
    \frac{\partial^2
    u}{\partial t^{2}} \label{eq:Dgl6}
\end{align}
Wir setzen nun \eqref{eq:Dgl2} und \eqref{eq:Dgl6} in \eqref{eq:Dgl3} und \eqref{eq:Dgl1} und \eqref{eq:Dgl5} in
\eqref{eq:Dgl4} ein und erhalten
\begin{align}
    \frac{\partial^{2} u}{\partial x^{2}} &= R^{\prime} G^{\prime} u(x,t) + (R^{\prime} C^{\prime} + L^{\prime}
    G^{\prime}) \frac{\partial u(x, t)}{\partial t} + L^{\prime} C^{\prime} \frac{\partial^{2} u(x,t)}{\partial t^{2}}
    \\[1.5ex]
    \frac{\partial^{2} i}{\partial x^{2}} &= R^{\prime} G^{\prime} i(x,t) + (R^{\prime} C^{\prime} + L^{\prime}
    G^{\prime}) \frac{\partial i(x, t)}{\partial t} + L^{\prime} C^{\prime} \frac{\partial^{2} i(x, t)}{\partial t^{2}}
\end{align}
Im Allgemeinen wird man diese partiellen Differentialgleichungen nur numerisch lösen können. Im Folgenden werden wir
einige Spezialfälle betrachten, die es uns erlauben, das qualitative Verhalten der Signalausbreitung zu beschreiben.


\subsection{Die Gleichungen der verlustlosen Leitung}
Wir betrachten zunächst den Fall, dass die auftretenden Verlustterme $R^{\prime}$ und $G^{\prime}$ vernachlässigbar
sind, Die Gleichungen \eqref{eq:Dgl1} und \eqref{eq:Dgl2} vereinfachen zu
\begin{align}
    \frac{\partial u}{\partial z} &= - L^{\prime} \frac{\partial i}{\partial t} \label{eq:Dgl7} \\
    \frac{\partial i}{\partial z} &= - C^{\prime} \frac{\partial u}{\partial t} \label{eq:Dgl8}
\end{align}






\chapter{Telegraphenleitungen}

\section{Verlustloses Model}

\section{Verlustbehaftetes Model}

\section{Abschlüsse}


%\bibliographystyle{plain} % oder ein anderer Stil, z.B. alpha, abbrv, unsrt
%\bibliography{Proseminar} % Name Ihrer .bib-Datei OHNE Endung .bib
\printbibliography

\end{document}